% theme: geometric_proofs_with_vectors
% 叩き台。平行・共線・垂直・重心・共点を、位置ベクトルと内積で証明する。
\MajorQuestionLesson{ベクトルによる図形の証明}
\SetRowSpace{4mm}

% figure: 凸四角形 ABCD と各辺の中点 P,Q,R,S
\TypeQuestion{凸四角形 $\mathrm{ABCD}$ の辺 $\mathrm{AB},\mathrm{BC},\mathrm{CD},\mathrm{DA}$ の中点をそれぞれ $\mathrm{P},\mathrm{Q},\mathrm{R},\mathrm{S}$ とする。四角形 $\mathrm{PQRS}$ が平行四辺形であることを、位置ベクトルを用いて証明しなさい。}{prove_varignon_parallelogram}
\QQ{証明しなさい。}{}

% figure: 三角形 ABC と中点 D,E,F
\TypeQuestion{$\Tri{ABC}$ の辺 $\mathrm{BC},\mathrm{CA},\mathrm{AB}$ の中点をそれぞれ $\mathrm{D},\mathrm{E},\mathrm{F}$ とする。$\Tri{ABC}$ と $\Tri{DEF}$ の重心が一致することを、位置ベクトルを用いて証明しなさい。}{prove_same_centroid_of_medial_triangle}
\QQ{証明しなさい。}{}

% figure: 三角形 ABC と3本の中線
\TypeQuestion{$\Tri{ABC}$ の3本の中線が1点で交わり、その交点が各中線を頂点の側から $2:1$ に内分することを、位置ベクトルを用いて証明しなさい。}{prove_concurrence_of_medians}
\QQ{証明しなさい。}{}

% figure: D∈BC, E∈CA, F∈AB。BD:DC=2:1, CE:EA=3:2, AF:FB=1:3
\TypeQuestion{$\Tri{ABC}$ の辺 $\mathrm{BC},\mathrm{CA},\mathrm{AB}$ 上にそれぞれ点 $\mathrm{D},\mathrm{E},\mathrm{F}$ をとり、$\Lseg{BD}:\Lseg{DC}=2:1$、$\Lseg{CE}:\Lseg{EA}=3:2$、$\Lseg{AF}:\Lseg{FB}=1:3$ とする。3直線 $\mathrm{AD},\mathrm{BE},\mathrm{CF}$ が1点で交わることを、位置ベクトルを用いて証明しなさい。}{prove_concurrence_of_cevians}
\QQ{証明しなさい。}{}

\LessonPageBreak
\SetRowSpace{4mm}

% figure: ひし形 ABCD と2本の対角線
\TypeQuestion{ひし形 $\mathrm{ABCD}$ の対角線 $\mathrm{AC}$ と $\mathrm{BD}$ が垂直であることを、ベクトルの内積を用いて証明しなさい。}{prove_rhombus_diagonals_perpendicular}
\QQ{証明しなさい。}{}

% figure: 正方形 ABCD、P は BC の中点、Q は CD の中点
\TypeQuestion{正方形 $\mathrm{ABCD}$ の辺 $\mathrm{BC},\mathrm{CD}$ の中点をそれぞれ $\mathrm{P},\mathrm{Q}$ とする。このとき、$\Lseg{AP}\perp\Lseg{BQ}$ であることを、ベクトルの内積を用いて証明しなさい。}{prove_perpendicular_lines_in_square}
\QQ{証明しなさい。}{}

% figure: 外心 O、重心 G、垂心 H
\TypeQuestion{正三角形ではない $\Tri{ABC}$ の外心を $\mathrm{O}$、重心を $\mathrm{G}$、垂心を $\mathrm{H}$ とする。3点 $\mathrm{O},\mathrm{G},\mathrm{H}$ が同一直線上にあり、$\Lseg{OG}:\Lseg{GH}=1:2$ となることを、位置ベクトルを用いて証明しなさい。}{prove_euler_line}
\QQ{証明しなさい。}{}

% figure: 平行四辺形 ABCD、P∈AB、Q∈AD、R=DP∩BQ
\TypeQuestion{平行四辺形 $\mathrm{ABCD}$ の辺 $\mathrm{AB},\mathrm{AD}$ をそれぞれ $2:1$ に内分する点を $\mathrm{P},\mathrm{Q}$ とし、線分 $\mathrm{DP}$ と線分 $\mathrm{BQ}$ の交点を $\mathrm{R}$ とする。3点 $\mathrm{A},\mathrm{R},\mathrm{C}$ が同一直線上にあることを、位置ベクトルを用いて証明しなさい。}{prove_collinearity_in_parallelogram}
\QQ{点 $\mathrm{R}$ の位置ベクトルを求め、3点 $\mathrm{A},\mathrm{R},\mathrm{C}$ が同一直線上にあることを示しなさい。}{}
\QQ{$\Lseg{AR}:\Lseg{RC}$ を求めなさい。}{}
